
\section{第一章 JavaScript 语言入门基础}
\subsection{1.1 语言定位与核心特性}
\begin{itemize}
\item JavaScript诞生与版本迭代发展
\item 弱类型动态脚本语言属性
\item 直译式解释执行运行特点
\item 浏览器客户端运行场景
\item Node.js服务端运行场景
\item 网页交互驱动设计思想
\item 单线程同步执行机制
\item 原型链面向对象特性
\end{itemize}

\subsection{1.2 运行环境与代码引入}
\begin{itemize}
\item 浏览器JS解析引擎作用
\item V8高性能引擎执行原理
\item Node.js后端运行环境搭建
\item script标签内嵌脚本编写
\item 外部独立JS文件引入加载
\item 事件触发延迟执行脚本
\end{itemize}

\subsection{1.3 基础语法书写规范}
\begin{itemize}
\item 标识符合法字符组成规则
\item 大小写敏感命名区分特性
\item 标识符驼峰命名行业规范
\item 系统关键字分类与作用
\item 保留字使用限制约束
\item 单行注释编写格式用途
\item 多行块注释编写格式用途
\item 语句分号结束分隔规则
\item 大括号代码块边界划分
\end{itemize}

\section{第二章 变量常量与完整数据类型}
\subsection{2.1 变量与常量声明方式}
\begin{itemize}
\item var老式变量声明与变量提升
\item let块级作用域变量定义规则
\item const常量声明不可修改特性
\item 变量初始化赋值操作
\item 变量二次赋值数值改写
\item 全局变量作用生效范围
\item 局部函数作用域范围
\item 块级作用域隔离访问规则
\end{itemize}

\subsection{2.2 七大原始基础数据类型}
\begin{itemize}
\item number数值整数浮点存储类型
\item string字符串文本不可变类型
\item boolean布尔真假逻辑判定类型
\item null空对象占位类型
\item undefined未定义默认空类型
\item symbol唯一不可重复标识类型
\item bigint超大整数高精度数值类型
\end{itemize}

\subsection{2.3 引用复合数据类型}
\begin{itemize}
\item object对象键值封装结构
\item array数组有序元素容器
\item function函数功能引用类型
\item 引用类型堆内存存储特点
\end{itemize}

\subsection{2.4 数据类型检测判定}
\begin{itemize}
\item typeof简易基础类型检测
\item Object.prototype精准类型判断
\item instanceof引用对象类型校验
\end{itemize}

\subsection{2.5 数据类型相互转换}
\begin{itemize}
\item 运算隐式自动类型转换规则
\item 字符串转数值强制转换
\item 数值转字符串格式转换
\item 任意数据转布尔判定转换
\item 转换过程数据丢失问题
\end{itemize}

\section{第三章 全类别运算符体系}
\subsection{3.1 算术运算符}
\begin{itemize}
\item 加法运算数值拼接双重特性
\item 减法运算数值计算规则
\item 乘法运算数值计算规则
\item 除法整数小数运算区分
\item 取模余数运算计算逻辑
\item 前置自增执行先后顺序
\item 后置自增执行先后顺序
\item 前置自减后置自减运算逻辑
\end{itemize}

\subsection{3.2 赋值与复合赋值运算符}
\begin{itemize}
\item 基础赋值数据覆盖逻辑
\item 加减乘除取模复合简写赋值
\end{itemize}

\subsection{3.3 关系比较运算符}
\begin{itemize}
\item 相等松散等值判断规则
\item 全等严格类型数值双判断
\item 不等与不全等反向判定
\item 大于小于大小边界比较
\item 大于等于小于等于判定
\item 关系运算布尔结果返回
\end{itemize}

\subsection{3.4 逻辑运算符}
\begin{itemize}
\item 逻辑与全真条件判定规则
\item 逻辑或一真条件判定规则
\item 逻辑非布尔值取反运算
\item 逻辑短路与截断执行特性
\item 逻辑短路或截断执行特性
\item 多条件逻辑组合嵌套判断
\end{itemize}

\subsection{3.5 位运算与特殊运算符}
\begin{itemize}
\item 二进制按位与或异或取反
\item 二进制左右位移运算逻辑
\item 三元条件判断运算符格式
\item 逗号多表达式顺序执行
\item 运算符优先级排序层级
\item 运算符左右结合执行顺序
\end{itemize}

\section{第四章 程序流程控制全部结构}
\subsection{4.1 顺序执行结构}
\begin{itemize}
\item 自上而下逐行代码执行逻辑
\item 顺序语句先后执行约束规则
\end{itemize}

\subsection{4.2 分支判断结构}
\begin{itemize}
\item if单分支条件判断语句
\item if-else双分支二选一逻辑
\item else if多分支区间逐级匹配
\item switch开关常量匹配分支
\item case分支匹配执行规则
\item break阻断分支穿透作用
\item default默认兜底分支逻辑
\item 分支语句多层嵌套使用方式
\end{itemize}

\subsection{4.3 多类型循环遍历结构}
\begin{itemize}
\item while先判断后执行条件循环
\item do-while先执行后判断循环
\item for三段式计数循环完整结构
\item for-in遍历对象自有属性
\item for-of遍历数组字符串元素
\end{itemize}

\subsection{4.4 循环跳转与嵌套控制}
\begin{itemize}
\item break终止跳出当前单层循环
\item continue跳过单次循环迭代
\item 双层多层循环嵌套运行顺序
\item 无限循环构成与终止条件设置
\end{itemize}

\section{第五章 数组与对象核心操作}
\subsection{5.1 数组基础定义与读写}
\begin{itemize}
\item 数组字面量构造函数创建
\item 数组下标从零起始索引规则
\item 通过下标存取修改数组元素
\item length属性获取数组长度
\end{itemize}

\subsection{5.2 数组增删改查基础方法}
\begin{itemize}
\item push尾部追加元素操作
\item pop尾部删除元素操作
\item unshift头部插入元素
\item shift头部移除元素
\item splice定点截取替换元素
\end{itemize}

\subsection{5.3 数组遍历与高阶处理方法}
\begin{itemize}
\item 普通for循环遍历数组
\item forEach遍历执行回调函数
\item filter条件筛选数组元素
\item map映射返回新数组数据
\item reduce累计聚合数组数值
\end{itemize}

\subsection{5.4 数组排序拼接截取操作}
\begin{itemize}
\item sort数组元素排序规则
\item reverse数组反转顺序
\item concat多数组合并拼接
\item slice截取指定区间子数组
\end{itemize}

\subsection{5.5 对象基础创建与访问}
\begin{itemize}
\item 对象字面量快速创建实例
\item new构造函数创建对象
\item 点语法访问对象属性
\item 中括号语法动态访问属性
\item 对象属性新增删除操作
\end{itemize}

\subsection{5.6 对象方法与属性遍历}
\begin{itemize}
\item 对象内部方法定义与调用
\item in运算符判定属性存在性
\item Object.keys获取属性名数组
\item 对象自有属性与继承属性区分
\end{itemize}

\section{第六章 函数编程与作用域闭包}
\subsection{6.1 函数多种定义形式}
\begin{itemize}
\item 函数声明式定义提升特性
\item 函数表达式定义写法
\item 箭头函数精简语法格式
\item 匿名函数无名称函数结构
\end{itemize}

\subsection{6.2 函数参数传递机制}
\begin{itemize}
\item 形式参数定义默认赋值
\item 实际参数传入匹配规则
\item arguments类数组参数集合
\item rest剩余参数收集多数据
\item 基本类型值传递拷贝特性
\item 引用类型地址传递特性
\end{itemize}

\subsection{6.3 函数返回值与调用}
\begin{itemize}
\item return语句终止函数并返回数据
\item 单值多数据返回处理方式
\item 无返回值默认undefined结果
\item 函数直接调用执行流程
\end{itemize}

\subsection{6.4 递归函数原理与使用}
\begin{itemize}
\item 函数自身嵌套调用逻辑
\item 递归递进执行过程
\item 递归回溯返回流程
\item 递归终止边界条件设置
\end{itemize}

\subsection{6.5 作用域层级查找规则}
\begin{itemize}
\item 全局作用域全局变量共享
\item 函数局部作用域隔离变量
\item 块级作用域变量访问限制
\item 作用域链逐层变量查找机制
\end{itemize}

\subsection{6.6 闭包概念与实际特性}
\begin{itemize}
\item 内层函数引用外层变量构成闭包
\item 闭包延长局部变量生命周期
\item 闭包私有数据缓存作用
\item 闭包内存占用泄漏问题
\end{itemize}

\section{第七章 DOM文档对象模型操作}
\subsection{7.1 DOM树形结构基础}
\begin{itemize}
\item 文档对象模型整体概念
\item 网页节点层级树形关系
\item 元素节点文本节点属性节点
\item 根节点文档顶层对象document
\end{itemize}

\subsection{7.2 多种DOM元素获取方式}
\begin{itemize}
\item 根据ID精准获取唯一元素
\item 根据类名获取元素集合
\item 根据标签名批量查找元素
\item CSS选择器单元素查询
\item CSS选择器多元素集合查询
\end{itemize}

\subsection{7.3 元素内容读写操作}
\begin{itemize}
\item innerHTML读写标签包含HTML结构
\item innerText纯文本内容读写
\item textContent标准文本内容操作
\end{itemize}

\subsection{7.4 元素属性动态管理}
\begin{itemize}
\item 获取元素固有属性数值
\item 设置修改元素固有属性
\item 自定义属性存取读取
\item 属性新增删除清除操作
\end{itemize}

\subsection{7.5 元素样式动态控制}
\begin{itemize}
\item 行内样式直接赋值修改
\item 类名切换控制样式效果
\item 获取元素实际计算样式
\end{itemize}

\subsection{7.6 DOM节点增删复制操作}
\begin{itemize}
\item 创建新空白元素节点
\item 末尾头部插入页面节点
\item 指定位置插入节点元素
\item 移除页面存在节点
\item 节点深浅拷贝复制操作
\end{itemize}

\section{第八章 网页事件处理机制}
\subsection{8.1 事件基础触发原理}
\begin{itemize}
\item 用户行为触发事件机制
\item 浏览器监听捕获事件
\item 常见鼠标键盘页面事件
\end{itemize}

\subsection{8.2 三种事件绑定方式}
\begin{itemize}
\item HTML行内标签绑定事件
\item DOM元素简单事件绑定
\item addEventListener标准事件监听
\end{itemize}

\subsection{8.3 事件对象属性方法}
\begin{itemize}
\item 事件对象event自动生成
\item 获取触发事件源元素
\item 鼠标坐标按键判定信息
\item 键盘按键编码识别获取
\end{itemize}

\subsection{8.4 事件传播三个阶段}
\begin{itemize}
\item 事件捕获自上而下传播
\item 目标阶段触发当前元素事件
\item 事件冒泡自下而上扩散
\end{itemize}

\subsection{8.5 事件阻止与委托}
\begin{itemize}
\item 阻止事件冒泡扩散行为
\item 阻止默认网页原生行为
\item 事件委托父级代理监听子元素
\item 事件委托性能优化原理
\end{itemize}

\section{第九章 内置常用核心对象}
\subsection{9.1 String字符串对象方法}
\begin{itemize}
\item 字符串字符位置查找匹配
\item 截取指定区间子字符串
\item 字符串替换指定内容
\item 字符串分割转为数组
\item 大小写格式统一转换
\item 首尾空格清除处理
\end{itemize}

\subsection{9.2 Math数学运算对象}
\begin{itemize}
\item 最大值最小值求取计算
\item 四舍五入向上向下取整
\item 平方开方基础数学运算
\item 随机小数整数范围生成
\end{itemize}

\subsection{9.3 Date日期时间对象}
\begin{itemize}
\item 时间对象当前时间创建
\item 年月日时分秒毫秒获取
\item 时间戳毫秒总数换算
\item 自定义时间初始化创建
\item 日期时间格式拼接转换
\end{itemize}

\section{第十章 BOM浏览器对象模型}
\subsection{10.1 window顶层窗口对象}
\begin{itemize}
\item 全局window对象顶层作用
\item 弹窗提示确认输入交互框
\item 一次性延时定时器使用
\item 循环间隔定时器执行
\item 定时器清除终止运行
\end{itemize}

\subsection{10.2 location地址栏对象}
\begin{itemize}
\item 网址协议域名路径参数解析
\item 页面跳转新地址访问
\item 页面原地刷新重载操作
\item URL查询参数获取拆分
\end{itemize}

\subsection{10.3 history浏览历史对象}
\begin{itemize}
\item 历史页面记录数量获取
\item 后退上一页浏览页面
\item 前进下一页浏览页面
\item 指定步数跳转历史记录
\end{itemize}

\subsection{10.4 navigator浏览器信息}
\begin{itemize}
\item 浏览器类型版本信息读取
\item 操作系统设备环境判定
\end{itemize}

\section{第十一章 异步编程体系}
\subsection{11.1 同步与异步执行差异}
\begin{itemize}
\item 同步阻塞等待执行逻辑
\item 异步非阻塞后台执行特性
\item 事件循环主线程任务队列
\end{itemize}

\subsection{11.2 回调函数异步模式}
\begin{itemize}
\item 回调函数触发执行时机
\item 多层回调嵌套地狱问题
\end{itemize}

\subsection{11.3 Promise异步对象}
\begin{itemize}
\item pending初始等待状态
\item resolve成功完成状态
\item reject失败异常状态
\item then成功回调链式调用
\item catch异常错误捕获处理
\item finally最终必定执行回调
\end{itemize}

\subsection{11.4 async与await语法糖}
\begin{itemize}
\item async修饰异步函数定义
\item await暂停等待异步结果
\item 同步写法实现异步逻辑
\item 异步代码异常捕获处理
\end{itemize}

\section{第十二章 AJAX网络请求与JSON}
\subsection{12.1 前后端交互与HTTP协议}
\begin{itemize}
\item 客户端服务端数据交互流程
\item HTTP请求响应基础模型
\item 请求头响应头信息作用
\item 状态码标识请求结果状态
\end{itemize}

\subsection{12.2 原生AJAX请求流程}
\begin{itemize}
\item XMLHttpRequest请求对象创建
\item 请求地址方式初始化配置
\item GET查询参数拼接携带
\item POST请求请求体数据发送
\item 监听请求状态变化回调
\item 响应报文数据解析处理
\end{itemize}

\subsection{12.3 JSON数据交换格式}
\begin{itemize}
\item JSON键值标准格式规范
\item JSON与JS对象语法区别
\item JSON字符串转JS对象解析
\item JS对象转为JSON序列化字符串
\end{itemize}

\section{第十三章 ES6及进阶新语法}
\subsection{13.1 变量解构赋值}
\begin{itemize}
\item 数组位置对应解构取值
\item 对象属性匹配解构赋值
\item 解构默认值容错设置
\end{itemize}

\subsection{13.2 拓展与剩余运算符}
\begin{itemize}
\item 拓展运算符展开数组对象
\item 数组合并对象合并简写
\item 剩余参数收拢多个实参
\end{itemize}

\subsection{13.3 对象语法简写优化}
\begin{itemize}
\item 属性名变量简写赋值
\item 函数方法简写定义格式
\item 对象动态属性名创建
\end{itemize}

\subsection{13.4 模块化导入导出}
\begin{itemize}
\item 默认导出暴露模块成员
\item 按需单独导出指定内容
\item 默认导入外部模块数据
\item 解构导入多项模块成员
\end{itemize}

\subsection{13.5 Set与Map新数据结构}
\begin{itemize}
\item Set自动去重集合特性
\item Map任意类型键值映射结构
\item 集合增删查改遍历操作
\end{itemize}
